% The FLFM light path seen from the camera (supervisor, 2026-11-23 meeting).
% The tube lens and Fourier lens are cut: the MLA is drawn in the back pupil
% of the objective. Adds the camera relay and the sensor with its pixels,
% which Figure flfm-path omits. Five lenslets drawn (N = 5); not to scale.
\begin{tikzpicture}[
    x=1cm,y=1cm,
    font=\footnotesize,
    ink/.style={text=black!85},
    lens/.style={<->,>={Stealth[length=2mm]},draw=black!75,line width=0.8pt},
    ray/.style={draw=blue!60!black,line width=0.7pt},
    relayray/.style={draw=blue!45,line width=0.5pt},
    slice/.style={draw=black!25,line width=0.4pt},
    plane/.style={draw=black!45,dashed,line width=0.5pt},
    subimg/.style={draw=orange!75!black,line width=2.2pt,opacity=0.55},
    dim/.style={<->,>={Stealth[length=1.6mm]},draw=black!70,line width=0.5pt},
    lead/.style={draw=black!55,line width=0.4pt},
    lab/.style={font=\footnotesize,ink,inner sep=1.5pt,align=center}
]
% ---- geometry (cm; not to scale) ----
% object 0 | objective 1.7 | back pupil = MLA 3.4 | MLA focal plane 5.4
% | relay lens 1 7.4 | relay Fourier plane 9.4 | relay lens 2 10.7 | sensor 12.0
\def\xo{1.7}  \def\xp{3.4}  \def\xf{5.4}
\def\xa{7.4}  \def\fa{2.0}  \def\fb{1.3}
\pgfmathsetmacro{\xb}{\xa+\fa+\fb}
\pgfmathsetmacro{\xs}{\xb+\fb}
\pgfmathsetmacro{\mc}{\fb/\fa}          % camera relay magnification (schematic)
\def\dm{0.56}                           % lenslet pitch d_MLA
\def\hp{1.4}                            % half of the pupil (5 pitches)
\def\sfr{0.8}                           % sub-image width as a fraction of the cell

% optical axis
\draw[draw=black!30,line width=0.4pt] (-0.4,0) -- (\xs+0.4,0);

% planes
\draw[plane] (0,-1.6) -- (0,1.6);
\draw[plane] (\xf,-1.6) -- (\xf,1.6);

% lenses
\draw[lens] (\xo,-1.6) -- (\xo,1.6);
\draw[lens] (\xa,-1.45) -- (\xa,1.45);
\draw[lens] (\xb,-1.0) -- (\xb,1.0);

% microlens array in the back pupil: 5 lenslets
\foreach \k in {-2,...,2} {
    \draw[draw=black!70,fill=black!6,line width=0.5pt]
        (\xp-0.07,{\k*\dm-0.5*\dm}) rectangle (\xp+0.07,{\k*\dm+0.5*\dm});
}

% rays of one on-axis emitter: objective collimates, lenslets focus
\draw[ray] (0,0) -- (\xo,\hp);
\draw[ray] (0,0) -- (\xo,-\hp);
\draw[ray] (\xo,\hp) -- (\xp,\hp);
\draw[ray] (\xo,-\hp) -- (\xp,-\hp);
\foreach \k in {-2,...,2} {
    \draw[slice] (\xo,{\k*\dm-0.5*\dm}) -- (\xp,{\k*\dm-0.5*\dm});
    \draw[ray] (\xp,{\k*\dm-0.5*\dm}) -- (\xf,{\k*\dm});
    \draw[ray] (\xp,{\k*\dm+0.5*\dm}) -- (\xf,{\k*\dm});
}
\draw[slice] (\xo,{2.5*\dm}) -- (\xp,{2.5*\dm});

% sub-images in the MLA focal plane: one per lenslet, narrower than a cell
\foreach \k in {-2,...,2} {
    \draw[subimg] (\xf,{\k*\dm-0.5*\sfr*\dm}) -- (\xf,{\k*\dm+0.5*\sfr*\dm});
    \fill[orange!80!black] (\xf,{\k*\dm}) circle (0.045);
}

% camera relay (4f): each spot is imaged, inverted, onto the sensor
\foreach \k in {-2,...,2} {
    \draw[relayray] (\xf,{\k*\dm}) -- (\xa,{\k*\dm}) -- (\xa+\fa,0)
        -- (\xb,{-\k*\dm*\mc}) -- (\xs,{-\k*\dm*\mc});
}

% sensor with pixel ticks
\draw[draw=black!70,fill=black!12,line width=0.5pt] (\xs,-1.15) rectangle (\xs+0.12,1.15);
\foreach \y in {-1.1,-1.045,...,1.12} {
    \draw[draw=black!45,line width=0.25pt] (\xs,\y) -- (\xs+0.12,\y);
}
\foreach \k in {-2,...,2} {
    \draw[subimg] (\xs-0.02,{-\k*\dm*\mc-0.5*\sfr*\dm*\mc})
        -- (\xs-0.02,{-\k*\dm*\mc+0.5*\sfr*\dm*\mc});
}

% dimensions at the array and the sensor
\draw[dim] (\xp-0.4,{1.5*\dm}) -- (\xp-0.4,{2.5*\dm});
\node[lab,anchor=east,fill=white,inner sep=1pt] at (\xp-0.45,{2.0*\dm}) {$d_{MLA}$};
\draw[dim] (\xs+0.3,-1.15) -- (\xs+0.3,1.15);
\node[lab,anchor=south,inner sep=1pt] at (\xs+0.35,1.17) {$D_{cam}$};
\draw[draw=black!70,line width=0.6pt] (\xs+0.75,{-0.5*\dm*\mc}) -- (\xs+0.75,{0.5*\dm*\mc});

% focal lengths along the bottom
\def\yb{-2.6}
\draw[dim] (0,\yb) -- (\xo,\yb);        \node[lab,fill=white] at ({0.5*\xo},\yb) {$f_{obj}$};
\draw[dim] (\xo,\yb) -- (\xp,\yb);      \node[lab,fill=white] at ({0.5*(\xo+\xp)},\yb) {$f_{obj}$};
\draw[dim] (\xp,\yb) -- (\xf,\yb);      \node[lab,fill=white] at ({0.5*(\xp+\xf)},\yb) {$f_{MLA}$};
\draw[dim] (\xf,\yb) -- (\xa,\yb);      \node[lab,fill=white] at ({0.5*(\xf+\xa)},\yb) {$f_{c1}$};
\draw[dim] (\xa,\yb) -- (\xa+\fa,\yb);  \node[lab,fill=white] at ({\xa+0.5*\fa},\yb) {$f_{c1}$};
\draw[dim] (\xa+\fa,\yb) -- (\xb,\yb);  \node[lab,fill=white] at ({\xa+\fa+0.5*\fb},\yb) {$f_{c2}$};
\draw[dim] (\xb,\yb) -- (\xs,\yb);      \node[lab,fill=white] at ({\xb+0.5*\fb},\yb) {$f_{c2}$};

% names
\node[lab,anchor=south] at (0,1.65) {object\\plane};
\node[lab,anchor=south] at (\xo,1.65) {objective};
\node[lab,anchor=south] at (\xp,1.65) {MLA in the\\back pupil};
\node[lab,anchor=south] at (\xf,1.65) {sub-images\\(MLA focal\\plane)};
\draw[draw=black!60,line width=0.5pt]
    (\xa-0.1,1.52) -- (\xa-0.1,1.62) -- (\xb+0.1,1.62) -- (\xb+0.1,1.52);
\node[lab,anchor=south] at ({0.5*(\xa+\xb)},1.72) {camera relay\\$M_{cam}=f_{c2}/f_{c1}$};
\node[lab,anchor=south] at (\xs+0.06,1.65) {sensor};
\node[lab,anchor=north,align=center] at (\xf+0.15,-1.65) {cell $\le d_{MLA}$:\\no overlap};

% ---- inset: the pixels of one view (zoom of the central cell) ----
\def\xi{13.4}   \def\px{0.2}
\pgfmathsetmacro{\xiE}{\xi+10*\px}
\draw[lead,dashed] (\xs+0.75,{0.5*\dm*\mc}) -- (\xi,1.0);
\draw[lead,dashed] (\xs+0.75,{-0.5*\dm*\mc}) -- (\xi,-1.0);
\draw[draw=black!25,line width=0.3pt] (\xi,-1.0) grid[step=\px] (\xiE,1.0);
\draw[draw=black!70,line width=0.5pt] (\xi,-1.0) rectangle (\xiE,1.0);
\fill[orange!80!black,opacity=0.55] ({\xi+5*\px},0) circle ({\px});
\fill[black!40] ({\xi+8*\px},{8*\px-1.0}) rectangle ({\xi+9*\px},{9*\px-1.0});
\node[lab,anchor=south] at ({0.5*(\xi+\xiE)},1.05) {one view\\on the sensor};
\draw[dim] (\xi,-1.2) -- (\xiE,-1.2);
\node[lab,anchor=north,align=center] at ({0.5*(\xi+\xiE)},-1.25)
    {$FOV$ on $\approx N_{px}/N$ pixels\\gray: one pixel $P$ ($p$ in object space)\\spot: $\Delta x_{LF}=S_r\,p$};
\end{tikzpicture}
